% Created 2025-11-05 Wed 10:09
% Intended LaTeX compiler: pdflatex
\documentclass[]{beamer}
			  \usepackage{mdframed}
			  \usepackage{tcolorbox}\usepackage{etoolbox}
			  \BeforeBeginEnvironment{minted}{\begin{tcolorbox}[title=Code]}\AfterEndEnvironment{minted}{\end{tcolorbox}}
			  \newenvironment{question}{\begin{tcolorbox}[colback=blue!5!white,colframe=blue!75!black,title={\bf Question }]}{\end{tcolorbox}}
\usepackage[utf8]{inputenc}
\usepackage[T1]{fontenc}
\usepackage{graphicx}
\usepackage{longtable}
\usepackage{wrapfig}
\usepackage{rotating}
\usepackage[normalem]{ulem}
\usepackage{amsmath}
\usepackage{amssymb}
\usepackage{capt-of}
\usepackage{hyperref}
\usepackage{tikz}
\usepackage[english]{babel}
\usepackage{minted} \usepackage[outputdir=.]{minted} \definecolor{CamBlue}{RGB}{163,193,173} \setbeamercolor{frametitle}{bg=CamBlue} \usepackage[orientation=landscape,size=custom,width=16,height=9,scale=0.5]{beamerposter} \graphicspath{{./bin/lectures/07/},{./bin/lectures/04/},{./bin/lectures/06/}} \usepackage{resizegather} \usepackage{cancel}
\usetheme{metropolis}
\author{Carl Henrik Ek}
\date{5th of November, 2025}
\title{Machine Learning and the Physical World}
\subtitle{Lecture 8 : Probabilistic Numerics}
\author{\texorpdfstring{Carl Henrik Ek - \url{che29@cam.ac.uk}}{Carl Henrik Ek}} \institute[INST]{\url{http://carlhenrik.com}} \titlegraphic{\includegraphics[height=0.7cm]{bin/cambridge_logo.png}} \usepackage{mathtools} \usepackage{graphicx} \mathtoolsset{showonlyrefs} \usepackage{tikz} \usetikzlibrary{arrows,shadows,shapes,backgrounds,decorations,snakes,fit,calc,timeline} \usepackage{tikz-3dplot} \usepackage{resizegather} \usepackage{xmpmulti} \usepackage{mathtools} \usefonttheme[onlymath]{serif} \PassOptionsToPackage{hyperref,x11names}{xcolor} \definecolor{electricblue}{HTML}{05ADF3} \usepackage{tocloft} \renewcommand{\cftsecleader}{\cftdotfill{\cftdotsep}} \usepackage[breaklinks=true,linktocpage,xetex]{hyperref} \hypersetup{colorlinks, citecolor=electricblue,filecolor=electricblue,linkcolor=electricblue,urlcolor=electricblue} \metroset{background=light,bloc=empty,progressbar=frametitle,block=empty} \usepackage[citestyle=authoryear-icomp, maxcitenames=1, maxbibnames=4, doi=false, isbn=false, eprint=false, bibstyle=authoryear, backend=bibtex, hyperref=true, url=false, natbib=true]{biblatex} \addbibresource{~/Dropbox/work/bibliography/references.bib}
\hypersetup{
 pdfauthor={Carl Henrik Ek},
 pdftitle={Machine Learning and the Physical World},
 pdfkeywords={},
 pdfsubject={},
 pdfcreator={Emacs 30.2 (Org mode 9.7.11)}, 
 pdflang={English}}
\begin{document}

\maketitle
\begin{frame}[label={sec:orgeeb76e8}]{}
\begin{center}
\includegraphics[height=0.7\textheight]{./bin/lectures/02/distribution.png}
\end{center}
\end{frame}
\begin{frame}[label={sec:org871db68}]{What do we do with uncertainty?}
\begin{center}
\includegraphics[height=0.7\textheight]{bin/lectures/07/headinthesand.png}
\end{center}
\end{frame}
\begin{frame}[label={sec:org9184c3c}]{Bayesian Optimisation}
\begin{center}
\includegraphics[width=.9\linewidth]{./bin/lectures/06/bo-aq-int-ei-4.png}
\end{center}
\end{frame}
\begin{frame}[label={sec:org4418ebf}]{}
\begin{center}
\includegraphics[height=0.7\textheight]{uncertaintyquantification02.png}
\end{center}
\[
y_i = f_i + \epsilon
\]
\end{frame}
\begin{frame}[label={sec:org4e329a5}]{}
\begin{quote} %% Poincare
"The need for probability only arises out of uncertainty: It has no place if we are certain that we know all aspects of a problem. \alert{But our lack of knowledge also must not be complete, otherwise we would have nothing to evaluate}. There is thus a spectrum of degrees of uncertainty. While the probability for the sixth decimal digit of a number in a table of logarithms to equal 6 is 1/10 a priori, in reality, all aspects of the corresponding problem are well determined, and, if we wanted to make the effort, we could find out its exact value. \alert{The same holds for interpolation, for the integration methods of Cotes or Gauss, etc}"

-- Henri Poincare, 1896
\end{quote}
\end{frame}
\begin{frame}[label={sec:org291e501}]{Uncertainty Quantification}
\begin{description}[<+->]
\item[{Aleatoric/Stochastic}] "Randomness" inherent in system, or noise in our measurement of system
\item[{Epistemic}] Uncertainty related to our ignorance of a the underlying system
\item[{Computational}] \emph{Uncertainty related to finite computation, or intractable computations}
\end{description}
\end{frame}
\begin{frame}[label={sec:org5865480}]{}
\begin{center}
Data + Model \(\overbrace{\rightarrow}^{Compute}\) Prediction
\end{center}
\end{frame}
\begin{frame}[label={sec:orga6da087}]{Computational Decisions}
\begin{itemize}[<+->]
\item Computation is expensive, how much knowledge will I gain from computing more?
\item What should I compute in order to reduce my uncertainty as much as possible?
\item How much should I trust the computation I have done?
\item How precise should I do down-stream tasks based on the information from a specific computation?
\end{itemize}
\end{frame}
\begin{frame}[label={sec:orgac0817f}]{Why Probabilistic Numerics?}
\begin{quote} %% Von Neumann
"[round-off errors] are strictly very complicated but uniquely defined number theoretical functions [of the inputs], yet our ignorance of their true nature is such that we best treat them as random variables."

-- \cite{neumann-1947-numer-inver}
\end{quote}
\end{frame}
\section{I believe in\ldots{}}
\label{sec:org922d249}
\begin{frame}[label={sec:org302ef37}]{}
\begin{center}
\includegraphics[width=.9\linewidth]{./bin/lectures/07/forrester.png}
\label{}
\end{center}
\end{frame}
\begin{frame}[label={sec:org0abf879}]{Formalisation [\cite{cockayne17_bayes_probab_numer_method}]}
\begin{center}
\includegraphics[height=0.6\textheight]{./bin/lectures/07/statisticallearning0.pdf}
\label{}
\end{center}

\[
F : p(\mathcal{D}) \rightarrow p(\mathcal{Y}\vert \mathcal{X})
\]
\end{frame}
\begin{frame}[label={sec:orgca737f5}]{Formalisation [\cite{cockayne17_bayes_probab_numer_method}]}
\begin{center}
\includegraphics[height=0.6\textheight]{./bin/lectures/07/statisticallearning0_1.pdf}
\label{}
\end{center}

\[
A\circ\mathcal{S} \approx F\circ p(\mathcal{D})
\]
\end{frame}
\begin{frame}[label=cockayne4]{Formalisation}
\begin{center}
\includegraphics[height=0.6\textheight]{./bin/lectures/07/statisticallearning5.pdf}
\label{}
\end{center}

\[
A\circ S\circ p(\mathcal{D}) \approx p(\mathcal{D})
\]
\end{frame}
\begin{frame}[label={sec:org224daaf}]{Formalisation [\cite{cockayne17_bayes_probab_numer_method}]}
\begin{center}
\includegraphics[height=0.6\textheight]{./bin/lectures/07/statisticallearning0_2.pdf}
\label{}
\end{center}

\[
A\circ S \circ p(\mathcal{D}) \approx F\circ p(\mathcal{D})
\]
\end{frame}
\begin{frame}[label={sec:org6d25258}]{Bayesian Optimisation}
\begin{center}
\includegraphics[width=.9\linewidth]{./bin/lectures/06/bo-aq-int-ei-4.png}
\end{center}
\end{frame}
\begin{frame}[label={sec:orgcc47e40}]{Bayesian Optimisation}
\begin{center}
\includegraphics[width=.9\linewidth]{./bin/lectures/07/statisticallearning_bo.pdf}
\label{}
\end{center}
\end{frame}
\begin{frame}[label={sec:org2edccdd}]{Formalisation [\cite{cockayne17_bayes_probab_numer_method}]}
\begin{center}
\includegraphics[height=0.6\textheight]{./bin/lectures/07/statisticallearning0_2.pdf}
\label{}
\end{center}

\[
A\circ S \circ p(\mathcal{D}) \approx F\circ p(\mathcal{D})
\]
\end{frame}
\begin{frame}[label={sec:org68615e5}]{Numerical Computations\footnote{\url{https://www.cs.toronto.edu/\~duvenaud/talks/odes\_runge\_kutta\_nips.pdf}}}
\begin{description}
\item[{Linear Algebra}] given \(As=y\) estimate \(x\) s.t. \(Ax=b\)
\item[{Optimisation}] given \(\nabla f(x_i)\) estimate \(x\) s.t. \(\nabla f(x) = 0\)
\item[{Analysis}] given \(f(x,t)\) estimate \(x(t)\) s.t. \(\textrm{d}x = f(x,t)\)
\item[{Quadrature}] given \(f(x_i)\) estimate \(\int_{a}^b f(x)\textrm{d}x\)
\end{description}
\end{frame}
\begin{frame}[label={sec:org32c69e7}]{Quantity of Interest}
\begin{center}
\includegraphics[width=.9\linewidth]{./bin/lectures/07/forrester_integrate.png}
\label{}
\end{center}
\end{frame}
\begin{frame}[label={sec:orgefe7608}]{}
\begin{quote} %% Integration
Integration is a significant numerical problem in many fields of science and engineering. It is a key step in inference, where it is encountered when averaging over the many states of the world consistent with observed data. Indeed, a provocative Bayesian view is that integration is the single challenge separating us from systems that fully automate statistics. More speculatively still, such systems may even exhibit artificial intelligence (ai).

-- Hennig, Osbourne, Kersting
\end{quote}
\end{frame}
\begin{frame}[label={sec:org9e46181}]{Integration}
\[
       F := \int f(x) \textrm{d}\nu(x)
       \]
\begin{itemize}
\item \(\nu(x)\) is the measure that we are integrating over
\end{itemize}
\end{frame}
\begin{frame}[label={sec:orgecef624}]{Integration}
\[
\underbrace{p(\mathcal{D})}_{F} = \int \underbrace{p(\mathcal{D}\mid \theta)}_{f(\theta)} \underbrace{p(\theta) \textrm{d}\theta}_{\textrm{d}\nu(\theta)}
\]
\begin{itemize}
\item marginalisation\footnote{think of computing the evidence} is integration over the prior probability measure on the parameter
\end{itemize}
\end{frame}
\begin{frame}[label={sec:org61397d4}]{Quadrature [\cite{cockayne17_bayes_probab_numer_method}]}
\begin{center}
\includegraphics[height=0.6\textheight]{./bin/lectures/07/statisticallearning7.pdf}
\label{}
\end{center}

\[
A\circ\mathcal{S} \approx \int f(x)\textrm{d}x
\]
\end{frame}
\begin{frame}[label={sec:orgc5f5c3b}]{Quadrature}
\begin{center}
\includegraphics[width=.9\linewidth]{./bin/lectures/07/quadrature0.png}
\label{}
\end{center}
\end{frame}
\begin{frame}[label={sec:org05c6810}]{Quadrature}
\begin{center}
\includegraphics[width=.9\linewidth]{./bin/lectures/07/quadrature1.png}
\label{}
\end{center}
\end{frame}
\begin{frame}[label={sec:org4132c1e}]{Numerical Computation}
\begin{quote} %% Definition
A numerical method \alert{estimates} a function's \alert{latent} property \alert{given} the result of computations.
\end{quote}
\begin{onlyenv}<2->
\(\frac{\text{Numerical algorithms}}{\text{Statistical inference}}\) takes data in the form of \(\frac{\text{evaluations of computations}}{\text{measurements of observed variables}}\) and aims to return predictions of the quantity of interest.
\end{onlyenv}
\begin{onlyenv}<3->
Should we think about computation as inference?
\end{onlyenv}
\end{frame}
\begin{frame}[label={sec:org0c406c8}]{Use of Computational Uncertainty}
\[
p(\hat{f} \mid S, A)
\]

\begin{description}[<+->]
\item[{Decision}] which algorithm to use when
\item[{Decision}] efficient use of expensive algorithms
\item[{Decision}] when to stop computation
\item[{Decision}] effect on downstream tasks
\end{description}
\end{frame}
\section{Bayesian Quadrature}
\label{sec:org243fe43}
\begin{frame}[label=indefiniteintegral]{Quantity of interest}
\[
F := \int_{-3}^3 \underbrace{e^{-(sin(3x))^2-x^2}}_{f(x)}\textrm{d}x
\]

\begin{itemize}
\item \(f(x)\) fully specified and deterministic
\item \(F\) is deterministic
\item \(F\) cannot be computed analytically
\end{itemize}
\end{frame}
\begin{frame}[label={sec:org5855cb7}]{Integration}
\begin{center}
\includegraphics[width=.9\linewidth]{./bin/lectures/07/integra.png}
\label{}
\end{center}
\end{frame}
\begin{frame}[label={sec:orgf4c7f8e}]{What we would like}
\[
p(F\mid Y)
\]
\begin{itemize}[<+->]
\item given that I have seen data \(Y\) what is my belief about the integral
\item allows for "active learning"
\item exploration/exploitation etc.
\end{itemize}
\end{frame}
\begin{frame}[label={sec:orge695878}]{Emulation}
\begin{center}
\includegraphics[width=.9\linewidth]{./bin/lectures/06/bo-aq-int-gp-15.png}
\end{center}
\end{frame}
\begin{frame}[label={sec:org8dfecda}]{Quadrature}
\begin{center}
\includegraphics[width=.9\linewidth]{./bin/lectures/07/bayesian-quadrature-intuition-00.png}
\label{}
\end{center}
\end{frame}
\begin{frame}[label={sec:org98df982}]{Quadrature}
\begin{center}
\includegraphics[width=.9\linewidth]{./bin/lectures/07/bayesian-quadrature-intuition-01.png}
\label{}
\end{center}
\end{frame}
\begin{frame}[label={sec:org16751b8}]{Quadrature}
\begin{center}
\includegraphics[width=.9\linewidth]{./bin/lectures/07/bayesian-quadrature-intuition-02.png}
\label{}
\end{center}
\end{frame}
\begin{frame}[label={sec:orgd7ec559}]{Quadrature}
\begin{center}
\includegraphics[width=.9\linewidth]{./bin/lectures/07/bayesian-quadrature-intuition-03.png}
\label{}
\end{center}
\end{frame}
\begin{frame}[label={sec:orgcdc6037}]{Quadrature}
\begin{center}
\includegraphics[width=.9\linewidth]{./bin/lectures/07/bayesian-quadrature-intuition-04.png}
\label{}
\end{center}
\end{frame}
\begin{frame}[label={sec:orgf45f7e9}]{Quadrature}
\begin{center}
\includegraphics[width=.9\linewidth]{./bin/lectures/07/bayesian-quadrature-intuition-05.png}
\label{}
\end{center}
\end{frame}
\begin{frame}[label={sec:org024ee61}]{Quadrature}
\begin{center}
\includegraphics[width=.9\linewidth]{./bin/lectures/07/bayesian-quadrature-intuition-06.png}
\label{}
\end{center}
\end{frame}
\begin{frame}[label={sec:org6f748c9}]{Model}
\[
F := \int_{-3}^3 \underbrace{e^{-(sin(3x))^2-x^2}}_{f(x)}\textrm{d}x
\]
\alert{Knowledge}
\begin{itemize}
\item \(f(x)\) strictly positive \(\Rightarrow\) \(F>0\)
\item bounded above by,
\[
  f(x) \leq e^{-x^2}
  \]
\item Therefore,
\[
  0< F < \int_{-\infty}^{\infty} e^{-x^2} \textrm{d}x = \sqrt{\pi}
  \]
\end{itemize}
\end{frame}
\begin{frame}[label=dataandknowledge]{}
\begin{center}
\includegraphics[width=.9\linewidth]{./bin/lectures/data11.png}
\label{}
\end{center}
\end{frame}
\begin{frame}[label={sec:org53ef3f3}]{Integration}
\[
F := \int f(x) \textrm{d}\nu(x)
\]
\begin{itemize}
\item \(\nu(x)\) is the measure that we are integrating over
\end{itemize}
\end{frame}
\begin{frame}[label={sec:org164b0da}]{Bayesian Quadrature [\cite{OHagan1991}]}
\begin{align}
p(F, Y) &= \int p(F\mid f)p(Y\mid f)p(f)\textrm{d}f
\only<2->{\\&=\int \delta\left(F-\int_{\mathcal{X}} f\textrm{d}x \right)\prod_{i}^N \delta(y_i - f(x_i))p(f)\textrm{d}f}
\end{align}
\end{frame}
\begin{frame}[label={sec:org912d120}]{Bayesian Quadrature [\cite{OHagan1991}]}
\[
p\left( \begin{array}{c} Y \\ F\end{array} \right) = \mathcal{N}\left(\left[\begin{array}{c}\mathbf{m}_X\\\int m_X(x)\textrm{d}x\end{array}\right],\left[\begin{array}{cc}k(X,X) & \int k(X,x)\textrm{d}x\\\int k(x,X)\textrm{d}x & \int\int k(x,x')\textrm{d}x\textrm{d}x'\end{array}\right]\right)
\]
\begin{itemize}[<+->]
\item We can derive \(p(F\mid Y)\) through our normal conditioning procedure
\item \(p(F\mid Y) = \mathcal{N}(\mu_F, k_F)\) is a uni-variate Gaussian
\end{itemize}
\end{frame}
\begin{frame}[label={sec:org88cdfaa}]{Bayesian Quadrature}
\[
p\left( \begin{array}{c} Y \\ F\end{array} \right) = \mathcal{N}\left(\left[\begin{array}{c}\mathbf{m}_X\\\int m_X(x)\textrm{d}x\end{array}\right],\left[\begin{array}{cc}k(X,X) & \int k(X,x)\textrm{d}x\\\int k(x,X)\textrm{d}x & \int\int k(x,x')\textrm{d}x\textrm{d}x'\end{array}\right]\right)
\]
\begin{itemize}
\item We can derive \(p(F\mid Y)\) through our normal conditioning procedure
\item \(p(F\mid Y) = \mathcal{N}(\mu_F, k_F)\) is a uni-variate Gaussian
\end{itemize}
\end{frame}
\begin{frame}[label={sec:org10026a7}]{Statistical Inference}
\begin{center}
\includegraphics[width=.9\linewidth]{./bin/lectures/bq-inference0.png}
\label{}
\end{center}
\end{frame}
\begin{frame}[label={sec:org45027d1}]{Statistical Inference}
\begin{center}
\includegraphics[width=.9\linewidth]{./bin/lectures/bq-inference1.png}
\label{}
\end{center}
\end{frame}
\begin{frame}[label={sec:org7f2aa25}]{Statistical Inference}
\begin{center}
\includegraphics[width=.9\linewidth]{./bin/lectures/bq-inference2.png}
\label{}
\end{center}
\end{frame}
\begin{frame}[label={sec:org8e41aab}]{Statistical Inference}
\begin{center}
\includegraphics[width=.9\linewidth]{./bin/lectures/bq-inference3.png}
\label{}
\end{center}
\end{frame}
\begin{frame}[label={sec:org62dfcb8}]{Bayesian Quadrature}
\[
p\left( \begin{array}{c} Y \\ F\end{array} \right) = \mathcal{N}\left(\left[\begin{array}{c}\mathbf{m}_X\\\int m_X(x)\textrm{d}x\end{array}\right],\left[\begin{array}{cc}k(X,X) & \int k(X,x)\textrm{d}x\\\int k(x,X)\textrm{d}x & \int\int k(x,x')\textrm{d}x\textrm{d}x'\end{array}\right]\right)
\]
\begin{itemize}
\item We can derive \(p(F\mid Y)\) through our normal conditioning procedure
\item \(p(F\mid Y) = \mathcal{N}(\mu_F, k_F)\) is a uni-variate Gaussian
\item \(p(Y\mid F) = \mathcal{N}(\mu_Y, k_Y)\) \alert{is a Gaussian process}
\end{itemize}
\end{frame}
\begin{frame}[label={sec:org7c6ff1a}]{Integral Constrained Samples}
\begin{center}
\includegraphics[height=0.7\textheight]{./bin/lectures/07/quadrature_eq1.png}
\label{}
\end{center}

\[
F=4.0
\]
\end{frame}
\begin{frame}[label={sec:orgf9a508a}]{Integral Constrained Samples}
\begin{center}
\includegraphics[height=0.7\textheight]{./bin/lectures/07/quadrature_eq2.png}
\label{}
\end{center}

\[
F=1.0
\]
\end{frame}
\begin{frame}[label={sec:orgf664945}]{Information Operator\footnote{sometimes called a "Design Rule"}}
\begin{center}
\includegraphics[height=0.45\textheight]{./bin/lectures/07/bayesian-quadrature-intuition-04.png}
\end{center}
\begin{description}
\item[{Integrand variance}] \(\alpha(x) = k(x,x)\)
\item[{Integral Variance Reduction }] \(\alpha(x) = k_F(X,X) - k_F(X,x)\)
\end{description}
\end{frame}
\begin{frame}[label={sec:org9611865}]{Choice of Covariance}
\begin{center}
\includegraphics[height=0.7\textheight]{./bin/lectures/07/quadrature-covariance-samples.png}
\label{}
\end{center}


\[
p(f) = \mathcal{GP}\left(\boldsymbol{0}, \theta^2(\operatorname{min}(x,x')-\kappa)\right)
\]
\end{frame}
\begin{frame}[label={sec:orge7fa9f9}]{Quadrature Rule}
\[
\mathbb{E}[F] = \mathbb{E}_{p(f\mid Y)} \left[\int f(x)\textrm{d}x \right] = \sum_{i=1}^{N-1}\frac{x_{i+1}-x_i}{2}(f(x_{i+1})+f(x_i))
\]
\begin{itemize}
\item <2-> This is the normal trapezoid rule!!!
\item <3-> The algorithm is now tied to \alert{our belief} in the function!!!!
\item <4-> We can do inference over where to sample!!!!!!!!
\end{itemize}
\end{frame}
\begin{frame}[label={sec:org68e718b}]{Trapedzoid Rule}
\begin{definition}[Trapedzoid Rule]\label{sec:org71b3b70}
The trapezoidal rule is the posterior mean estimate for the integral \(F = \int_{a}^{b} f(x)\textrm{d}x\) under any centred Wiener process prior \(p(f) = \mathcal{GP}(0,k)\) with \(k(x, x′) = \theta^2(\operatorname{min}(x,x')-\kappa)\)
\end{definition}
\end{frame}
\begin{frame}[label={sec:org929b7d6}]{}
\begin{center}
\includegraphics[height=0.7\textheight]{bin/lectures/07/notimpressed.jpg}
\end{center}
\end{frame}
\begin{frame}[label=scientificprinciple]{The Scientific Principle}
\begin{center}
\includegraphics[height=0.7\textheight]{./bin/lectures/02/scientificprinciple.pdf}
\label{}
\end{center}
\begin{center}
Data + Model \(\overbrace{\rightarrow}^{Compute}\) Prediction
\end{center}
\end{frame}
\begin{frame}[label={sec:orgba271df}]{Trapedzoid Rule}
In calculus, the trapezoidal rule is a technique for numerical integration, i.e. approximating the definite integral:
\[
\int_{a}^{b}f(x)\textrm{d}x
\]
The trapezoidal rule works by approximating the region under the graph of the function \(f(x)\) as a trapezoid and calculating its area. This is easily calculated by noting that the area of the region is made up of a rectangle with width \(b-a\)  and height \(f(a)\) and a triangle of width \(b-a\) and height \(f(b)-f(a)\)
\end{frame}
\begin{frame}[label={sec:orgc09faf6}]{Least Squares Regression}
\begin{center}
\includegraphics[height=0.6\textheight]{./bin/lectures/07/linreg_posterior10.png}
\label{}
\end{center}


\begin{description}
\item[{Legendre (1805)}] algorithm that reduces "error"
\item[{Gauss (1809)}] statistical model assuming i.i.d. Gaussian noise
\end{description}
\end{frame}
\begin{frame}[label={sec:org2bec142}]{Factor Analysis}
\begin{center}
\includegraphics[height=0.5\textheight]{./bin/lectures/07/fa.png}
\label{}
\end{center}

\begin{description}
\item[{Spearman (1904)}] proposed an algorithm to extract "factors" from data \cite{Spearman:1904tx}
\item[{Hotelling (1936)}] concept of factor \alert{is} clearly defined through a statistical model \cite{Hotelling:1933ki}
\end{description}
\end{frame}
\begin{frame}[label={sec:orgc262929},fragile]{Why Probabilistic Numerics? \texttt{scipy.optimize.minimize}}
 \begin{minted}[]{python}
def minimize(fun, x0, args=(), method=None,
             jac=None, hess=None,
             hessp=None, bounds=None,
             constraints=(), tol=None,
             callback=None, options=None):
\end{minted}
\begin{description}
\item[{\alert{method}}] \texttt{Nelder-Mead}, \texttt{Powell}, \texttt{CG}, \texttt{BFGS}, \texttt{Newton-CG}, \texttt{L-BFGS-B}, \texttt{TNC} , \texttt{COBYLA} , \texttt{SLSQP} , \texttt{trust-constr} , \texttt{dogleg} , \texttt{trust-ncg} , \texttt{trust-exact} , \texttt{trust-krylov}
\end{description}
\end{frame}
\begin{frame}[label={sec:orgd452766}]{Probabilistic Numerics [\cite{Hennig:2015jf}]}
\begin{itemize}[<+->]
\item There are tons of numerical algorithms for every problem under the sun
\item They work really well
\item They give different results on the same problem
\item \emph{what is the prior they implement?}
\end{itemize}
\end{frame}
\begin{frame}[label={sec:org8b1fbbb}]{\href{https://parameterfree.com/2020/12/06/neural-network-maybe-evolved-to-make-adam-the-best-optimizer/}{Neural Networks (Maybe) Evolved to Make Adam The Best Optimizer}}
\emph{In a talk, Olivier Bousquet has described the deep learning community as a giant genetic algorithm: Researchers in this community are exploring the space of all variants of algorithms and architectures in a semi-random way. Things that consistently work in large experiments are kept, the ones not working are discarded. \ldots{}.
\ldots{}\ldots{}\ldots{}\ldots{}\ldots{}..
the community is evolving only one set of parameters (architectures, initialization strategies, hyperparameters search algorithms, etc.) keeping most of the time the optimizer fixed to Adam.}
\end{frame}
\begin{frame}[label={sec:org6915faf}]{"The Machine Learning Principle"\footnote{\cite{Chomsky:1980ud}}}
\begin{quotation} %% Noam Chomsky
"There is a notion of success \ldots{} which I think is novel in the history of science. It interprets success as approximating unanalyzed data."

-- Prof. Noam Chomsky
\end{quotation}
\end{frame}
\begin{frame}[label={sec:orgcc15372}]{}
\begin{center}
\includegraphics[height=0.7\textheight]{./bin/lectures/07/slicedbread.png}
\end{center}
\end{frame}
\begin{frame}[label={sec:org1c9cb8a}]{Assumptions: Algorithms}
\begin{columns}
\begin{column}{0.32\columnwidth}
\begin{center}
\includegraphics[width=.9\linewidth]{bin/lectures/02/eyepatch.png}
\end{center}
\end{column}
\begin{column}{0.32\columnwidth}
\begin{center}
\includegraphics[width=.9\linewidth]{bin/lectures/02/clownshoes.jpeg}
\end{center}
\end{column}
\begin{column}{0.32\columnwidth}
\begin{center}
\includegraphics[width=.9\linewidth]{bin/lectures/02/ministrysillywalks3.png}
\end{center}
\end{column}
\end{columns}
\begin{block}{Statistical Learning}
\[
\textcolor{red}{\mathcal{A}}_{\mathcal{H}}(\mathcal{S})
\]
\end{block}
\end{frame}
\section{Summary}
\label{sec:orgb461d5e}
\begin{frame}[label={sec:org5f4f809}]{Summary}
\begin{itemize}[<+->]
\item Probabilistic Numerics extends the notion of statistical inference to \alert{computation}\footnote{these thoughts have been around for a long time}
\item Computation is the process of extracting a latent property, machine learning is the statistical process of updating beliefs about latent properties
\item Computation is often not "truth" therefore we should quantify our ignorance about its results
\item Why?
\begin{itemize}
\item efficiency
\item down-stream tasks, uncertainty in computation should be part of decision
\item learning/understanding algorithms in relation to problems/data
\end{itemize}
\end{itemize}
\end{frame}
\begin{frame}[label={sec:org22a6010}]{When computation was expensive}
\begin{description}[<+->]
\item[{Albert Valentionvic Suldin (1924-1996)}] worked on error minimising estimators for numerical algorithms, how to \alert{design} algorithms from a statistical perspective
\item[{Frederick Michael Larkin (1936-1982)}] incorporating the notion of \alert{prior} knowledge into numerical algorithms to make robust calculations
\end{description}
\end{frame}
\begin{frame}[label={sec:org31b7472}]{}
\begin{center}
\includegraphics[height=0.7\textheight]{bin/lectures/07/pn-book.jpg}
\end{center}

\begin{center}
\url{http://probnumschool.org}
\end{center}
\end{frame}
\begin{frame}[label={sec:orge5c9dd1}]{}
\begin{center}
\alert{eof}
\end{center}
\end{frame}
\begin{frame}[label={sec:org93a39a4}]{Is BO PN?}
\begin{center}
\includegraphics[height=0.5\textheight]{./bin/lectures/07/statisticallearning8.pdf}
\label{}
\end{center}


\begin{description}[<+->]
\item[{Yes}] it uses a probablistic model as a proxy for decision loop
\item[{No}] the probabilistic model is not over the quantity of interest
\end{description}
\end{frame}
\begin{frame}[label=cockayne3]{Formalisation}
\begin{center}
\includegraphics[height=0.6\textheight]{./bin/lectures/07/statisticallearning2.pdf}
\label{}
\end{center}

\[
A\circ\mathcal{S}=\tilde{p}(\mathcal{D}) \approx p(\mathcal{D})
\]
\end{frame}
\begin{frame}[allowframebreaks]{}
 \printbibliography
\end{frame}
\end{document}
